% Options for packages loaded elsewhere
\PassOptionsToPackage{unicode}{hyperref}
\PassOptionsToPackage{hyphens}{url}
\documentclass[
  a4paper,
]{ltjsarticle}
\usepackage{xcolor}
\usepackage[margin=25mm]{geometry}
\usepackage{amsmath,amssymb}
\setcounter{secnumdepth}{-\maxdimen} % remove section numbering
\usepackage{iftex}
\ifPDFTeX
  \usepackage[T1]{fontenc}
  \usepackage[utf8]{inputenc}
  \usepackage{textcomp} % provide euro and other symbols
\else % if luatex or xetex
  \usepackage{unicode-math} % this also loads fontspec
  \defaultfontfeatures{Scale=MatchLowercase}
  \defaultfontfeatures[\rmfamily]{Ligatures=TeX,Scale=1}
\fi
\usepackage{lmodern}
\ifPDFTeX\else
  % xetex/luatex font selection
\fi
% Use upquote if available, for straight quotes in verbatim environments
\IfFileExists{upquote.sty}{\usepackage{upquote}}{}
\IfFileExists{microtype.sty}{% use microtype if available
  \usepackage[]{microtype}
  \UseMicrotypeSet[protrusion]{basicmath} % disable protrusion for tt fonts
}{}
\makeatletter
\@ifundefined{KOMAClassName}{% if non-KOMA class
  \IfFileExists{parskip.sty}{%
    \usepackage{parskip}
  }{% else
    \setlength{\parindent}{0pt}
    \setlength{\parskip}{6pt plus 2pt minus 1pt}}
}{% if KOMA class
  \KOMAoptions{parskip=half}}
\makeatother
\usepackage{longtable,booktabs,array}
\newcounter{none} % for unnumbered tables
\usepackage{calc} % for calculating minipage widths
% Correct order of tables after \paragraph or \subparagraph
\usepackage{etoolbox}
\makeatletter
\patchcmd\longtable{\par}{\if@noskipsec\mbox{}\fi\par}{}{}
\makeatother
% Allow footnotes in longtable head/foot
\IfFileExists{footnotehyper.sty}{\usepackage{footnotehyper}}{\usepackage{footnote}}
\makesavenoteenv{longtable}
\usepackage{graphicx}
\makeatletter
\newsavebox\pandoc@box
\newcommand*\pandocbounded[1]{% scales image to fit in text height/width
  \sbox\pandoc@box{#1}%
  \Gscale@div\@tempa{\textheight}{\dimexpr\ht\pandoc@box+\dp\pandoc@box\relax}%
  \Gscale@div\@tempb{\linewidth}{\wd\pandoc@box}%
  \ifdim\@tempb\p@<\@tempa\p@\let\@tempa\@tempb\fi% select the smaller of both
  \ifdim\@tempa\p@<\p@\scalebox{\@tempa}{\usebox\pandoc@box}%
  \else\usebox{\pandoc@box}%
  \fi%
}
% Set default figure placement to htbp
\def\fps@figure{htbp}
\makeatother
\setlength{\emergencystretch}{3em} % prevent overfull lines
\providecommand{\tightlist}{%
  \setlength{\itemsep}{0pt}\setlength{\parskip}{0pt}}
% Glyph map for lualatex (newunicodechar). Maps symbols that may lack font glyphs.
\usepackage{newunicodechar}
\usepackage{amssymb}
\newunicodechar{∎}{\ensuremath{\blacksquare}}
\newunicodechar{✓}{\ensuremath{\checkmark}}
\newunicodechar{✗}{\ensuremath{\times}}
\newunicodechar{⟹}{\ensuremath{\Longrightarrow}}
\newunicodechar{⟺}{\ensuremath{\Longleftrightarrow}}
\newunicodechar{↪}{\ensuremath{\hookrightarrow}}
\newunicodechar{⌊}{\ensuremath{\lfloor}}
\newunicodechar{⌋}{\ensuremath{\rfloor}}
\newunicodechar{≃}{\ensuremath{\simeq}}
\newunicodechar{≈}{\ensuremath{\approx}}
\newunicodechar{≤}{\ensuremath{\leq}}
\newunicodechar{≥}{\ensuremath{\geq}}
\newunicodechar{≠}{\ensuremath{\neq}}
\newunicodechar{→}{\ensuremath{\rightarrow}}
\newunicodechar{←}{\ensuremath{\leftarrow}}
\newunicodechar{↔}{\ensuremath{\leftrightarrow}}
\newunicodechar{⊥}{\ensuremath{\perp}}
\newunicodechar{√}{\ensuremath{\surd}}
\newunicodechar{∑}{\ensuremath{\sum}}
\newunicodechar{∏}{\ensuremath{\prod}}
\newunicodechar{∞}{\ensuremath{\infty}}
\newunicodechar{∈}{\ensuremath{\in}}
\newunicodechar{∀}{\ensuremath{\forall}}
\newunicodechar{∣}{\ensuremath{\mid}}
\newunicodechar{γ}{\ensuremath{\gamma}}
\newunicodechar{λ}{\ensuremath{\lambda}}
\newunicodechar{μ}{\ensuremath{\mu}}
\newunicodechar{ρ}{\ensuremath{\rho}}
\newunicodechar{σ}{\ensuremath{\sigma}}
\newunicodechar{φ}{\ensuremath{\varphi}}
\newunicodechar{ψ}{\ensuremath{\psi}}
\newunicodechar{θ}{\ensuremath{\theta}}
\newunicodechar{Δ}{\ensuremath{\Delta}}
\newunicodechar{Π}{\ensuremath{\Pi}}
\newunicodechar{Λ}{\ensuremath{\Lambda}}
\newunicodechar{τ}{\ensuremath{\tau}}
\newunicodechar{ε}{\ensuremath{\varepsilon}}
\newunicodechar{ω}{\ensuremath{\omega}}
\newunicodechar{π}{\ensuremath{\pi}}
\newunicodechar{ν}{\ensuremath{\nu}}
\newunicodechar{±}{\ensuremath{\pm}}
\newunicodechar{×}{\ensuremath{\times}}
\newunicodechar{·}{\ensuremath{\cdot}}
\usepackage{bookmark}
\IfFileExists{xurl.sty}{\usepackage{xurl}}{} % add URL line breaks if available
\urlstyle{same}
\hypersetup{
  pdftitle={90度高対称床の解析的厳密解------Z4位相ラベルが定める整数反対称生成子の \textbackslash sigma\_\{\textbackslash max\} 固有モードと閉形式スペクトル},
  pdfauthor={Noriaki Kihara (WF System Co., Ltd.), ORCID 0009-0004-6753-4020},
  hidelinks,
  pdfcreator={LaTeX via pandoc}}

\title{90度高対称床の解析的厳密解------Z4位相ラベルが定める整数反対称生成子の
\(\sigma_{\max}\) 固有モードと閉形式スペクトル}
\author{Noriaki Kihara (WF System Co., Ltd.), ORCID 0009-0004-6753-4020}
\date{2026-09-12}

\begin{document}
\maketitle

\textbf{シリーズ:} 自己無撞着な関係波閉鎖系の基礎（論文3／全10本）\\
\textbf{著者:} 木原 範昭（WF System Co., Ltd.）　\textbf{日付:}
2026-09-12\\
\textbf{Version DOI:} 10.5281/zenodo.22728996\\
\textbf{Concept DOI:} 10.5281/zenodo.22728995\\
\textbf{ORCID:} 0009-0004-6753-4020\\
\textbf{Zenodo:} https://zenodo.org/records/22728996\\

\begin{center}\rule{0.5\linewidth}{0.5pt}\end{center}

\subsection{要旨}\label{ux8981ux65e8}

論文2で点火の資格を満たすと確定した90度高対称床を、反復近似（make\_parent、\texttt{iters=1200},
\texttt{tol=1e-12}）ではなく\textbf{整数固有値問題として解析的に厳密に}構成し、その閉形式を\textbf{任意
\(N\)} に対して導出する。床の位相は大域位相を除き
\(\{0,90,180,270\}^\circ\) に量子化されるので各辺に Z4 ラベル
\(k_e\in\{0,1,2,3\}\) を与えると、生成子
\(K_{ef}=A_{ef}\sin(90^\circ(k_f-k_e))\) は\textbf{整数行列（成分
\(0,\pm1\)）}になる。決定論的種 \(k_e=(i+j)\bmod4\)
について、生成子を決めるのは mod 4 ではなく \textbf{parity}
\(p_e=(i+j)\bmod2\) だけで、論文1の実対称行列
\(W_{ef}=A_{ef}\sin^2\psi_{ef}\) が二部グラフ
\(W_{ef}=A_{ef}\mathbf 1_{p_e\neq p_f}\)（\(L(K_N)\)
を頂点番号和の偶奇で二分）になる。したがって床は
\(k_e=(i+j)\bmod4\to W\to\) Perron ベクトル
\(r>0\to z=Dr\)（\(D=\operatorname{diag}(e^{i90^\circ k_e})\)、論文1
\(D^{-1}KD=iW\)）で\textbf{一発で構成でき、離散ラベル反復は本質的に不要}である。床は
\(iK\) の\textbf{絶対値最大の負固有値} \(\mu=-\sigma_{\max}\)
の固有ベクトルで、位相は \(d\varphi/d\tau=-\mu=\sigma_{\max}>0\)
で前進し位相方程式と整合する。二部構造の縮約 Perron 問題（偶数 \(N\) は
\(2\times2\)、奇数 \(N\) は \(3\times3\)）から \(\sigma^2\)（偶数
\(N(N-2)\)／奇数 \(N^2-2N-1\)）・各辺振幅²（parity
クラスごとに有理数）・全スペクトル（\(BB^{\mathsf T}\)）が任意 \(N\)
で閉形式に出る（\textbf{sympy 記号で \(N\) 一般に検証、数値で
\(N=3\)--\(40\) 一致}）。さらに一般 \(N\) で非零スペクトルは \(O(N)\)
個であり、偶数 \(N\ge6\)・奇数 \(N\ge7\) では厳密に \(2(N-1)\)
個＝\(\operatorname{rank}K=2(N-1)\)・核次元
\((N-1)(N-4)/2\)（\(N=3,4,5\) は退化例）：\(M=N(N-1)/2=O(N^2)\)
の関係波空間のうち活動するのは \(O(N)\) 次元だけである。全
\(N=3,\ldots,40\)（38床）で one\_step 相対平衡残差
\(\sim10^{-15}\)--\(1.2\times10^{-14}\)（make\_parent の
\(\sim10^{-13}\)
を上回る）、位相は厳密に軸上、\(\|z\|^2=1\)。\textbf{零閉塞
\(\sum z^2=0\)
は床に課した追加条件ではなく、二部固有モード（\(\|x\|^2=\|y\|^2\)）であることから自動的に従う一般定理}（Perron
に限らず全非零固有モードで成立）で、二部 \(W\) の連結性も任意 \(N\ge3\)
で証明する（論文1の Perron
適用条件を回収）。床は「\(\sigma=\sqrt{\text{整数}}\)、各辺振幅
\(=\sqrt{\text{有理数}}\)（parity クラスで決まる）、位相
\(\in\{0,90,180,270\}\)」で完全に代数的に書き下せる。分類：解析的厳密解（parity
二部構造＋縮約 Perron 問題）＋記号（任意
\(N\)）/数値（\(N=3\)--\(40\)）検証。

\begin{center}\rule{0.5\linewidth}{0.5pt}\end{center}

\subsection{1.
課題------近似床を厳密化する}\label{ux8ab2ux984cux8fd1ux4f3cux5e8aux3092ux53b3ux5bc6ux5316ux3059ux308b}

論文2で用いた make\_parent
床は自己無撞着固有モードを反復で近似した解で、残差は
\(\sim10^{-13}\)。反復上限・許容カット・damping
で止めた近似をそのままコピーしても意味がない。90度床は位相が離散（Z4）に量子化されるという構造を持つので、これを使えば\textbf{整数固有値問題として代数的に厳密}に解ける。本論文はその厳密解を全
\(N\)（奇数を含む）で構成し、閉形式を確定する。

\subsection{2.
原理------Z4ラベルから整数反対称生成子へ}\label{ux539fux7406z4ux30e9ux30d9ux30ebux304bux3089ux6574ux6570ux53cdux5bfeux79f0ux751fux6210ux5b50ux3078}

床の位相を大域位相を除き4象限に量子化し、各辺に \(k_e\in\{0,1,2,3\}\)
を割り当てる。位相差は \(90^\circ(k_f-k_e)\) だから

\[K_{ef}=A_{ef}\sin\!\big(90^\circ(k_f-k_e)\big),\qquad (k_f-k_e)\bmod4:\ 0\to0,\ 1\to+1,\ 2\to0,\ 3\to-1,\]

すなわち \(K\) は成分 \(0,\pm1\) の\textbf{整数反対称行列}。床 \(z\)
はその \(\sigma_{\max}\) 固有モード、すなわち \(iK\)
の\textbf{絶対値最大の負固有値} \(\mu=-\sigma_{\max}\)
の固有ベクトル（\(z=\sigma_{\max}u-iKu\) の形、\(u\) は
\(K^2u=-\sigma_{\max}^2u\) の実固有ベクトル）。\(iK\)
はエルミートで固有値は共役対 \(\pm\sigma_j\)、\(\sigma_{\max}\)
が最大特異値（スペクトル半径）。床は位相が前進する分枝
\(\mu=-\sigma_{\max}\) を取り、\(d\varphi/d\tau=-\mu=\sigma_{\max}>0\)
が位相方程式 \(r\dot\varphi=\sum r_f\sin^2\Delta\varphi\ge0\)
と整合する（実データ・コードで確認：N=6 で make\_parent 床
\(\mu=-4.868\)、整数K床
\(\mu=-4.899\)、ともに最小固有値）。\(K^{\mathsf T}K\)
の固有値の平方根が定めるのは\textbf{特異値
\(\sigma=\sqrt{\text{整数}}\)} であって、各辺振幅 \(r_e\)
はそれに属する固有ベクトル（＝\(W\) の Perron
ベクトル）の成分であり、本構成では
\(r_e=\sqrt{\text{有理数}}\)（§5）の閉形式になる。

なお、複素自己無撞着問題 \(Kz=i\omega z\) が実対称非負整数行列
\(W_{ef}=A_{ef}\sin^2\psi_{ef}\) の固有値問題 \(Wr=\omega r\)
と厳密同値であること（\(D^{-1}KD=iW\)、\(D=\operatorname{diag}(e^{i\varphi_e})\)）と、\(W\)
既約下での Perron--Frobenius による正振幅解の存在・一意性は\textbf{論文1
が一般に与える}。本論文はその \(W\)＝整数 \(K\)
の具体形で、\(\sigma=\rho(W)\) と各辺振幅を閉形式（§5）に解く（\(W\) の
Perron 固有ベクトルと整数 \(K\) の \(\sigma_{\max}\)
モードは同一の床。\(r>0\)・主固有ベクトル一致を全 \(N\) で確認）。

\textbf{parity による二部構造。} \(W\) を決めるのは mod 4 ではなく
\textbf{mod 2}（parity）\(p_e=(i+j)\bmod2\)
だけである。\(\sin^2[90^\circ(k_f-k_e)]\) は \(k_f-k_e\)
が奇（\(p_e\neq p_f\)）なら \(1\)、偶（\(p_e=p_f\)）なら \(0\) だから

\[W_{ef}=A_{ef}\,\mathbf 1_{p_e\neq p_f}.\]

すなわち \(W\) は \(L(K_N)\) の頂点（＝\(K_N\) の辺
\(\{i,j\}\)）を頂点番号和の偶奇 \(p_e\)
で二分した\textbf{二部グラフの隣接行列}である。この \(W\) は任意
\(N\ge3\) で\textbf{連結（既約）}である（証明は §8：偶頂点集合
\(U\)（\(|U|=\lceil N/2\rceil\ge2\)）・奇頂点集合 \(V\)
に対し、異偶間辺どうしは内部辺を橋に、同偶奇辺は異偶間辺で接続する）。ゆえ
Perron ベクトル \(r>0\)
が一意（論文1）で、\(z=Dr\)（\(D=\operatorname{diag}(e^{i90^\circ k_e})\)）が
\(Kz=i\sigma_{\max}z\) を厳密に満たす。したがって床は種ラベルから
\(k\to W\to r\to z=Dr\)
で一発構成でき、離散ラベル反復（§3）は床の構成には不要である（反復床の
\(W\) が種 parity の \(W\) と機械精度一致することを全 \(N=3\)--\(40\)
で確認、‖差‖
\(\le3\times10^{-27}\)、\texttt{parity二部構造\_解析的証明\_20260912/}）。この二部構造から
§5・§6 の閉形式が任意 \(N\) で出る。

\textbf{二乗閉塞 \(\sum_e z_e^2=0\)
は二部構造の帰結（追加条件でない）。} 二部
\(W=\begin{pmatrix}0&B\\B^{\mathsf T}&0\end{pmatrix}\) の非零固有モード
\(r=(x,y)\) は \(By=\sigma x,\ B^{\mathsf T}x=\sigma y\) より
\(\sigma\|x\|^2=x^{\mathsf T}By=(B^{\mathsf T}x)^{\mathsf T}y=\sigma\|y\|^2\)、\(\sigma\neq0\)
ゆえ \(\|x\|^2=\|y\|^2\)（規格化で各 \(1/2\)）。一方
\(z_e=e^{i90^\circ k_e}r_e\) より
\(z_e^2=(-1)^{k_e}r_e^2=(-1)^{p_e}r_e^2\)（実数）だから

\[\sum_e z_e^2=\sum_{p_e=0}r_e^2-\sum_{p_e=1}r_e^2=\|x\|^2-\|y\|^2=0.\]

すなわち零閉塞 \(\sum z^2=0\)
は床に課した追加条件ではなく、\textbf{二部固有モードであることから自動的に生じる}（Perron
床に限らず二部 \(W\)
の任意の非零固有モードで成立。実部・交差項虚部が別々にゼロなのは
\(z_e^2\) が実であることの帰結）。全 \(N=3\)--\(40\)
の全非零固有モードで
\(|\sum z^2|\le10^{-15}\)・\(|\|x\|^2-\|y\|^2|\le10^{-15}\) を確認。

\subsection{3. 床の構成------種ラベルから Perron
で一発（反復は確認のみ）}\label{ux5e8aux306eux69cbux6210ux7a2eux30e9ux30d9ux30ebux304bux3089-perron-ux3067ux4e00ux767aux53cdux5fa9ux306fux78baux8a8dux306eux307f}

§2 より床は種ラベル \(k_e=(i+j)\bmod4\)（辺
\(e=\{i,j\}\)、\(0\le i<j\le N-1\)）から\textbf{反復なしで直接}構成できる：

\begin{enumerate}
\def\labelenumi{\arabic{enumi}.}
\tightlist
\item
  parity \(p_e=(i+j)\bmod2\) で二部行列
  \(W_{ef}=A_{ef}\mathbf 1_{p_e\neq p_f}\) を組む。
\item
  \(W\) の Perron ベクトル \(r>0\) と \(\sigma_{\max}=\rho(W)\)
  を取る（\(W\) 連結、論文1）。
\item
  \(D=\operatorname{diag}(e^{i90^\circ k_e})\) として床 \(z=Dr\)。論文1
  \(D^{-1}KD=iW\) より \(Kz=i\sigma_{\max}z\)（相対平衡）。
\end{enumerate}

この直接構成 \(z=Dr\) が one\_step 相対平衡残差 \(\le1.2\times10^{-14}\)
を満たすことを全 \(N=3\)--\(40\) で確認した。従来の生成器
\texttt{make\_90deg\_floor\_analytic\_v1.py}
は種から離散ラベル反復（\(iK\) の絶対値最大固有値 \(\mu=-\sigma_{\max}\)
の固有ベクトルを4象限へ再量子化→大域位相4通りで前ラベルに最一致→固定点/gauge等価サイクルで停止・最小残差床を採用）で床を求めるが、\textbf{反復床の
\(W\) は種 parity の \(W\) と機械精度一致}（‖差‖
\(\le3\times10^{-27}\)）し、床は種直接構成と\textbf{符号ゲージ}（\(z_e\to-z_e\)、\(k_e\to k_e+2\)、\(K\to SKS\)
で \(Kz=i\sigma_{\max}z\)・\(W\) 不変）を除いて同一である（ゲージ整合後
‖差‖
\(\le8\times10^{-15}\)）。したがって反復は床の構成には不要で、種が固定点（符号ゲージを除く）であることの確認に相当する。adjacency/one\_step
は正本
\texttt{run\_N3\_N40\_stage123\_v1.py}（SHA照合）から抽出（物理無変更）。各床の検算（\texttt{build\_analytic\_floor}）：\(\|z\|^2\)、\(|\sum z|\)、\(\Re(a^{\mathsf T}a-b^{\mathsf T}b)\)、\(a^{\mathsf T}b\)、\(|\sum z^2|\)、one\_step
相対平衡残差、軸上判定、象限数、振幅種数、\(\sigma\) を記録。

\subsection{\texorpdfstring{4. 結果（全
\(N=3,\ldots,40\)・38床）}{4. 結果（全 N=3,\textbackslash ldots,40・38床）}}\label{ux7d50ux679cux5168-n3ldots4038ux5e8a}

全床とも位相ちょうど軸上（\(0/90/180/270\)）、閉塞
\(\sum z^2=0\)（実部・交差項虚部とも機械ゼロ。§2
の二部構造の帰結）、\(\|z\|^2=1\)、相対平衡残差
\(\sim10^{-15}\)--\(1.2\times10^{-14}\)。

\begin{itemize}
\tightlist
\item
  \textbf{\(N=3\)}：3象限（T型）、\(|\sum z|=0.707\)、\(\sigma=\sqrt2=1.414\)。
\item
  \textbf{\(4\mid N\)（4,8,\ldots,40）}：重心
  \(|\sum z|\sim10^{-15}\)（厳密ゼロ）、振幅2値。
\item
  \textbf{\(N\equiv2\ (\mathrm{mod}\,4)\)（6,10,\ldots,38）}：重心
  \(0.236\to0.037\)（\(N\)増で減少）、振幅2値。
\item
  \textbf{奇数 \(N\)（5,7,\ldots,39）}：重心
  \(0.154\to0.019\)（減少）、振幅3値。\textbf{奇数も全て床}（軸上・閉塞ゼロ）。
\end{itemize}

重心の厳密ゼロは \(4\mid N\)（象限の完全 \(\pm\)
対称が可能）に限られ、他は parity 不均衡と Perron
解から決まる残差重心である（最小化ではなく決定量。閉形式は
§8。厳密ゼロは取れない）。精度比較（同一 \(N\)）：

{\def\LTcaptype{none} % do not increment counter
\begin{longtable}[]{@{}lll@{}}
\toprule\noalign{}
\(N\) & make\_parent 残差 & 整数K解析解 残差 \\
\midrule\noalign{}
\endhead
\bottomrule\noalign{}
\endlastfoot
6 & \(2.3\times10^{-13}\) & \(5.4\times10^{-16}\) \\
7 & \(3.5\times10^{-14}\) & \(1.7\times10^{-15}\) \\
8 & \(8.2\times10^{-14}\) & \(5.7\times10^{-16}\) \\
9 & \(8.2\times10^{-14}\) & \(1.0\times10^{-15}\) \\
\end{longtable}
}

全38床を複素平面グリッド図（図1）に収録。

\begin{figure}
\centering
\pandocbounded{\includegraphics[keepaspectratio,alt={図1: 90度理論床（整数K σ\_max 解析解, 全N=3..40）。位相 0/90/180/270・厳密相対平衡（残差\textasciitilde1e-15）・4\textbar N は重心厳密ゼロ。}]{figs/p3_analytic_floor_ja_fig1.png}}
\caption{図1: 90度理論床（整数K σ\_max 解析解, 全N=3..40）。位相
0/90/180/270・厳密相対平衡（残差\textasciitilde1e-15）・4\textbar N
は重心厳密ゼロ。}
\end{figure}

\textbf{図1} 90度理論床（整数K \(\sigma_{\max}\)
解析解、\(N=3\)--\(40\)）。

\subsection{\texorpdfstring{5. 閉形式（parity 縮約 Perron 問題から任意
\(N\)）}{5. 閉形式（parity 縮約 Perron 問題から任意 N）}}\label{ux9589ux5f62ux5f0fparity-ux7e2eux7d04-perron-ux554fux984cux304bux3089ux4efbux610f-n}

§2 の二部構造から Perron ベクトル \(r\) は parity
クラスで一定値をとる。各クラスの隣接本数を数えると縮約 Perron
問題になり、\(\sigma\) と各辺振幅が\textbf{任意 \(N\)}
で閉形式に解ける（sympy
記号検証、\texttt{parity二部構造\_解析的証明\_20260912/results/実行ログ}）。整数
\(K\) は反対称ゆえ \(K^2\) は整数対称（固有値 \(-\sigma^2\)）で
\(\sigma=\sqrt{\text{整数}}\)、各辺振幅² は有理数。

\textbf{偶数 \(N=2m\)（2クラス）。} 異偶間辺（\(p=1\)、\(i+j\) 奇、本数
\(N^2/4\)）と同偶奇内辺（\(p=0\)、\(i+j\) 偶、本数
\(N(N-2)/4\)）。各異偶間辺は同偶奇辺を \(N-2\)
本、各同偶奇辺は異偶間辺を \(N\) 本隣接に持つので、クラス代表値
\(x\)（異偶間）・\(y\)（同偶奇）は

\[\sigma x=(N-2)y,\qquad \sigma y=Nx\ \Longrightarrow\ \sigma^2=N(N-2),\]

\(\|z\|^2=1\) で規格化して \(x^2=2/N^2\)、\(y^2=2/(N(N-2))\)。

\textbf{奇数 \(N\)（3クラス）。} \(a=\tfrac{N+1}2\)
偶頂点・\(b=\tfrac{N-1}2\)
奇頂点。異偶間（\(ab=\tfrac{N^2-1}4\)）・偶側内部（\(\binom a2=\tfrac{N^2-1}8\)）・奇側内部（\(\binom b2=\tfrac{(N-1)(N-3)}8\)）。代表値
\(x\)（異偶間）・\(y\)（偶側内部）・\(z\)（奇側内部）は

\[\sigma x=\tfrac{N-1}2y+\tfrac{N-3}2z,\quad \sigma y=(N-1)x,\quad \sigma z=(N+1)x\ \Longrightarrow\ \sigma^2=N^2-2N-1,\]

規格化して
\(x^2=\tfrac{2}{(N-1)(N+1)}\)、\(y^2=\tfrac{2(N-1)}{(N+1)(N^2-2N-1)}\)、\(z^2=\tfrac{2(N+1)}{(N-1)(N^2-2N-1)}\)。

\textbf{\(\sigma^2\)}：偶数 \(N(N-2)\)／奇数 \(N^2-2N-1\)。数列
\(N=3\)--\(20\)：\(2,8,14,24,34,48,62,80,98,120,142,168,194,224,254,288,322,360\)。

\textbf{振幅²（規格化 \(\|z\|^2=1\)、値 × 本数）------クラス割当は
parity で確定する。}

偶数 \(N\)（2クラス・各クラス総パワー \(1/2\)）：

{\def\LTcaptype{none} % do not increment counter
\begin{longtable}[]{@{}lll@{}}
\toprule\noalign{}
振幅² & クラス（辺） & 本数 \\
\midrule\noalign{}
\endhead
\bottomrule\noalign{}
\endlastfoot
\(2/N^2\) & 異偶間（\(i+j\) 奇） & \((N/2)^2\) \\
\(2/(N(N-2))\) & 同偶奇内（\(i+j\) 偶） & \(N(N-2)/4\) \\
\end{longtable}
}

奇数 \(N\)（3クラス。\(N=3\) は奇側内部0本＝T型）：

{\def\LTcaptype{none} % do not increment counter
\begin{longtable}[]{@{}
  >{\raggedright\arraybackslash}p{(\linewidth - 4\tabcolsep) * \real{0.3333}}
  >{\raggedright\arraybackslash}p{(\linewidth - 4\tabcolsep) * \real{0.3333}}
  >{\raggedright\arraybackslash}p{(\linewidth - 4\tabcolsep) * \real{0.3333}}@{}}
\toprule\noalign{}
\begin{minipage}[b]{\linewidth}\raggedright
振幅²
\end{minipage} & \begin{minipage}[b]{\linewidth}\raggedright
クラス（辺）
\end{minipage} & \begin{minipage}[b]{\linewidth}\raggedright
本数
\end{minipage} \\
\midrule\noalign{}
\endhead
\bottomrule\noalign{}
\endlastfoot
\(2/((N-1)(N+1))\) & 異偶間（\(i+j\) 奇） & \((N^2-1)/4\) \\
\(2(N-1)/((N+1)(N^2-2N-1))\) & 偶側内部（\(i,j\) とも偶） &
\((N^2-1)/8\) \\
\(2(N+1)/((N-1)(N^2-2N-1))\) & 奇側内部（\(i,j\) とも奇） &
\((N-1)(N-3)/8\) \\
\end{longtable}
}

本数総和はいずれも \(M=N(N-1)/2\)。閉形式は sympy 記号（任意
\(N\)）で検証し、数値で \(N=3\)--\(40\) の各辺振幅² と一致（誤差
\(\le1.2\times10^{-16}\)）。偶数 \(N\) で重心が厳密ゼロなのは2クラスが各
\(1/2\) パワーで \(\pm\) 完全相殺するため。奇数 \(N\) は3クラスで
\(\pm\) 相殺が不完全なため小さな残差重心が残る（閉形式は §8）。

\subsection{\texorpdfstring{6.
固有振動スペクトルと活動次元（\(O(N)\)）}{6. 固有振動スペクトルと活動次元（O(N)）}}\label{ux56faux6709ux632fux52d5ux30b9ux30daux30afux30c8ux30ebux3068ux6d3bux52d5ux6b21ux5143on}

二部構造 \(W=\begin{pmatrix}0&B\\B^{\mathsf T}&0\end{pmatrix}\)
の非零固有値は
\(\pm\sqrt{\operatorname{eig}(BB^{\mathsf T})}\)。異偶間辺を（偶頂点)×(奇頂点)
の \(a\times b\) 格子で添字づけると

\[BB^{\mathsf T}=(a-2)I_a\otimes J_b+(b-2)J_a\otimes I_b+2J_a\otimes J_b,\]

固有値は
\(4ab-2(a+b)\)（\(\times1\)）、\(b(a-2)\)（\(\times(a-1)\)）、\(a(b-2)\)（\(\times(b-1)\)）、残り
\(0\)。\(a,b\) を代入して（偶数 \(N\) は \(a=b=N/2\)、奇数 \(N\) は
\(a=\tfrac{N+1}2,b=\tfrac{N-1}2\)）：

\begin{itemize}
\tightlist
\item
  \textbf{偶数
  \(N\ge6\)}：\(\lambda_{\max}^2=N(N-2)\)（多重度1）、\(\lambda_1^2=N(N-4)/4\)（多重度
  \(N-2\)）。
\item
  \textbf{奇数
  \(N\ge7\)}：\(\lambda_{\max}^2=N^2-2N-1\)（多重度1）、\(\lambda_1^2=(N-5)(N+1)/4\)（多重度
  \((N-3)/2\)）、\(\lambda_2^2=(N-1)(N-3)/4\)（多重度 \((N-1)/2\)）。
\end{itemize}

数値で全 \(N\) 一致（偶 \(N\ge6\)／奇 \(N\ge7\)）。\(N=3,4,5\)
は退化で一部の固有値族が消える（\(N=4\)：\(\lambda_1^2=0\)、\(N=5\)：\((N-5)(N+1)/4=0\)）。

\textbf{活動次元 \(O(N)\)。} 非零固有値は \(a+b-1=N-1\) 個（\(\pm\)
ペアで \(2(N-1)\)）。よって

\[\operatorname{rank}K=\operatorname{rank}W=2(N-1),\qquad \dim\ker=M-2(N-1)=\frac{(N-1)(N-4)}2\quad(\text{偶}N\ge6,\ \text{奇}N\ge7).\]

\(M=N(N-1)/2=O(N^2)\) の関係波空間のうち、90度床の生成子で活動するのは
\(O(N)\) 次元だけで、残りは \(\ker K\) に落ちる。小 \(N\)
は退化で例外：\(N=3\)（rank2, 核1）、\(N=4\)（rank2,
核4）、\(N=5\)（rank6, 核4）。

単純整数倍音系（基音の整数倍）になるのは \(N=3,4\) のみ。\(N\ge5\)
は一般に平方根比。生の固有値 \(\lambda\) を \(N\) 横断で比較すると
\(2\)--\(8\) 倍の機械精度整数比ペアが65組、位相刻みを込めた
\(q=\lambda/N\)
では11組検出される（\texttt{固有振動数解析\_90度系列\_N3\_N40\_20260910/}）。これは床が単一
carrier
を共巻きする単一固有モードであること（論文4：因数分解しても床上の周波数比は1:1）と整合する。

\subsection{7.
非90度理論床との関係}\label{ux975e90ux5ea6ux7406ux8ad6ux5e8aux3068ux306eux95a2ux4fc2}

{\def\LTcaptype{none} % do not increment counter
\begin{longtable}[]{@{}
  >{\raggedright\arraybackslash}p{(\linewidth - 8\tabcolsep) * \real{0.2000}}
  >{\raggedright\arraybackslash}p{(\linewidth - 8\tabcolsep) * \real{0.2000}}
  >{\raggedright\arraybackslash}p{(\linewidth - 8\tabcolsep) * \real{0.2000}}
  >{\raggedright\arraybackslash}p{(\linewidth - 8\tabcolsep) * \real{0.2000}}
  >{\raggedright\arraybackslash}p{(\linewidth - 8\tabcolsep) * \real{0.2000}}@{}}
\toprule\noalign{}
\begin{minipage}[b]{\linewidth}\raggedright
\end{minipage} & \begin{minipage}[b]{\linewidth}\raggedright
位相
\end{minipage} & \begin{minipage}[b]{\linewidth}\raggedright
生成子
\end{minipage} & \begin{minipage}[b]{\linewidth}\raggedright
解法
\end{minipage} & \begin{minipage}[b]{\linewidth}\raggedright
対称性
\end{minipage} \\
\midrule\noalign{}
\endhead
\bottomrule\noalign{}
\endlastfoot
巡回Fourier床（論文2の(C)） & \(2\pi(i+j)/N\)（N値） & 連続 &
距離クラス縮約固有ベクトル & 巡回対称（振幅は距離クラス一定） \\
90度整数K床（本論文） & \(0/90/180/270\)（4値） & \textbf{整数行列} &
\(\mu=-\sigma_{\max}\) 固有ベクトル & 距離クラス対称を破る（振幅は
parity クラスごと） \\
\end{longtable}
}

90度整数K床は距離クラス対称ではない（振幅は parity クラスごと。偶数
\(N\) で2値、奇数 \(N\) で3値）ので距離クラス縮約では作れず、二部 \(W\)
の Perron モード（＝整数 \(K\) の \(\sigma_{\max}\)
固有モード）として解くのが正しい。両床は別メンバーだが、ともに論文2の点火資格（不安定な相対平衡）を満たす。

\subsection{8.
主張分類・未決}\label{ux4e3bux5f35ux5206ux985eux672aux6c7a}

\begin{itemize}
\tightlist
\item
  分類：解析的厳密解（parity 二部構造＋縮約 Perron 問題）＋記号（任意
  \(N\), sympy）/数値（\(N=3\)--\(40\)）検証。判定：保持。
\item
  \textbf{解決済み（本版で閉形式化・一般証明）}：(1)
  振幅クラスの辺割当は parity クラスそのもの（異偶間＝\(i+j\)
  奇／偶側内部＝\(i,j\) とも偶／奇側内部＝\(i,j\)
  とも奇）に確定（§5）。(2) 奇数 \(N\)
  の残差重心も閉形式（\(D=N^2-2N-1\)）：
\end{itemize}

\[|\textstyle\sum z|=\begin{cases}0&N\equiv0\ (\mathrm{mod}\,4)\\[1mm] \sqrt2/N&N\equiv2\ (\mathrm{mod}\,4)\\[1mm] \sqrt{(N-1)/(2(N+1)D)}&N\equiv1\ (\mathrm{mod}\,4)\\[1mm] \sqrt{(N+1)/(2(N-1)D)}&N\equiv3\ (\mathrm{mod}\,4)\end{cases}\]

全 \(N=3\)--\(40\) で実測と一致（誤差 \(\le10^{-9}\)。例
\(N=5\)：0.154303、\(N=6\)：0.235702、\(N=7\)：0.140028）。(3) 二乗閉塞
\(\sum z^2=0\) は二部固有モードの帰結で、Perron
に限らず全非零固有モードで成立する一般定理（§2）。(4) \textbf{二部 \(W\)
の連結性を任意 \(N\ge3\) で一般証明（論文1の未決を回収）}：偶頂点
\(U\)（\(|U|=\lceil N/2\rceil\ge2\)）・奇頂点 \(V\)。異偶間辺
\((u,v),(u',v')\) は、\(u\neq u'\) なら
\((u,v)\)--\((u,u')\)--\((u',v')\)（\((u,u')\)
は偶側内部）、\(u=u'\)（\(v\neq v'\)、\(|V|\ge2\)）なら
\((u,v)\)--\((v,v')\)--\((u,v')\) で連結。同偶奇辺 \((u,u')\) は
\((u,u')\)--\((u,v)\)（任意 \(v\in V\)）でその成分に接続。ゆえ \(N\ge3\)
で連結＝既約、Perron--Frobenius が任意 \(N\ge3\) で適用可能（連結成分数
\(=1\) を \(N=3\)--\(40\) で数値確認）。(5)
符号ゲージ族の数え上げ：各辺独立に
\(k_e\to k_e+2t_e\)（\(S=\operatorname{diag}((-1)^{t_e})\)、\(K\to SKS\)・\(W\)
不変）で \(2^M\) 通り、全辺同時反転 \(S=-I\)（大域位相
\(\pi\)）を同一視して相対符号ゲージ \(2^{M-1}\)。 -
未決：本論文の範囲で構造的な未決は残らない。頂点置換など符号ゲージ以外の同値による床の分類（幾何的モジュライの整理）は本論文の対象外。

\subsection{関連研究（構造的一致・導出元でない）}\label{ux95a2ux9023ux7814ux7a76ux69cbux9020ux7684ux4e00ux81f4ux5c0eux51faux5143ux3067ux306aux3044}

\begin{itemize}
\tightlist
\item
  辺グラフのスペクトル（Cvetković et al.,
  代数的グラフ理論）：\(A=L(K_N)\) とその上の整数反対称作用素の固有値。
\item
  相対平衡（Marsden \& Ratiu）：床＝剛体回転相対平衡。
\end{itemize}

\subsection{参考文献}\label{ux53c2ux8003ux6587ux732e}

\begin{enumerate}
\def\labelenumi{\arabic{enumi}.}
\tightlist
\item
  木原範昭「自己無撞着な関係波閉鎖系におけるインフレーション的急拡大の機構」Concept
  DOI 10.5281/zenodo.22112008.
\item
  D. Cvetković, P. Rowlinson, S. Simić, \emph{An Introduction to the
  Theory of Graph Spectra}, Cambridge (2010).
\end{enumerate}

\subsection{再現}\label{ux518dux73fe}

生成器・検証・データは
\texttt{90度理論床\_整数K解析解\_全N\_20260909/}（\texttt{make\_90deg\_floor\_analytic\_v1.py}＝生成器、\texttt{closed\_form\_90deg\_floor\_verify\_v1.py}＝閉形式
sympy
厳密検証、\texttt{parents\_90deg\_floor\_analytic/parent\_90deg\_floor\_N\{00003..00040\}\_analytic.npz}＝38床、\texttt{summary\_90deg\_floor\_closed\_form.csv}、\texttt{manifest\_90deg\_floor\_analytic.json}、\texttt{fig\_90deg\_floor\_complex\_planes\_analytic.png}、\texttt{分析\_90度理論床\_整数K解析解\_全N\_20260909.md}、\texttt{run\_all.sh}、\texttt{SHA256SUMS.txt}）。parity
二部構造による任意 \(N\) 解析化（§2・§3・§5・§6・§8）の検証は
\texttt{同/parity二部構造\_解析的証明\_20260912/}（\texttt{verify\_parity\_bipartite\_analytic\_20260912.py}＝正本
adjacency/one\_step を SHA 照合 import の読出しのみ検証＋sympy
記号、\texttt{results/parity\_bipartite\_summary.csv}、\texttt{results/実行ログ\_20260912.log}、\texttt{分析\_parity二部構造\_解析的一本化\_20260912.md}、\texttt{README.md}、\texttt{run\_all.sh}、\texttt{SHA256SUMS.txt}）。固有振動スペクトルは
\texttt{同/固有振動数解析\_90度系列\_N3\_N40\_20260910/}。複素自己無撞着
≡ 実対称 \(W\) の同値（\(D^{-1}KD=iW\)）と Perron--Frobenius
は論文1。adjacency/one\_step は正本
\texttt{run\_N3\_N40\_stage123\_v1.py}（SHA256
\texttt{1abf2353fee2e4f56f05e7a6f149fd086885136beb61ab571b48a56b09691567}）から抽出・物理無変更。

\end{document}
